% New proof bridge. All earlier sources and audit files are preserved.
\subsection{A retained-phase bridge through Proposition 5.2}
\label{sec:lichtman-prop52-retained-phase-bridge-20261009}

\paragraph{Scope and the two explicit analytic inputs.}
This note supplies the small-factor, coefficient and error-composition
argument replacing the invocation of L23 Lemma~4.6 inside its
Proposition~5.2. It does not assume that the disputed phase-deletion
estimate is true. The following two inputs are stated separately,
so that a local replacement cannot silently stand in for a complete
large-residue dispersion theorem.

\begin{description}
\item[Input D.] The large-residue version of the preliminary dispersion
reduction, with $E=1$: for rough $\alpha_n,\Lambda_{d,r}$, full
auxiliary Siegel--Walfisz for $\alpha$, and divisor-bounded coefficients,
if $N>QD\log^{C}x$ and $M>\log^{C}x$, then
\[
 |\mathcal S|\ll_{A,B}
 x\log^{-A}x+M(DQ)^{1/2}\log^{O_B(1)}x
       (|\mathcal E_1|^{1/2}+|\mathcal E_2|^{1/2}).       \tag{D}
\]
The actual $\mathcal E_i$ are those at L23 primary text
lines1660--1682, including the factor $1/d$ in $\mathcal E_2$.
This input must include a valid treatment of $mn=a$, signs and
shifted-divisor moments at large $a$; simply citing the fixed-residue
Maynard proof is not a verification of that extension.

\item[Input S.] The bounded-coefficient well-factorable exponential-sum
estimate of Lemma~5.1, with fixed $d=x^{o(1)}$, its stated product
endpoints, and the bounded-$h$/independent-coefficient extensions
specified below. With source margin $\sigma$, its three size
conditions $7\sigma,14\sigma,14\sigma$ and $Q<Nx^{-\sigma}$
give $N^2Q^{-1}x^{-\sigma}$. The proof and spectral input, including
all occurrences of $a$, are audited separately. The extension does
not authorize arbitrary weights on the outer $q$ variable.
\end{description}

The separate $\mathcal E_1$ estimate used below is L23 Lemma~4.7,
whose statement has $\sup_{a\ne0}$ and whose asserted proof uses
Cauchy--Schwarz and the Weil bound, not Lemma~4.6. A failure of
Input D is therefore an earlier obstruction than the replacement
proved here. Conversely, once D and S are supplied with the stated
uniformity, the following derivation has no additional unproved
``precisely of the same form'' step.

\paragraph{Proposition (precise conditional bridge).}
Let $0<\sigma<1/100$, $0\le\kappa\le\sigma/100$ and
$0\le\theta<1/2$. Choose the scale $x=NM$; replacing an asymptotically
equivalent external scale by $NM$ only changes fixed constants and
can be accommodated by decreasing the available strict margins.
Let $a\in\mathbb Z\setminus\{0\}$ satisfy $|a|\le x^{1+\kappa}$.
Assume
\begin{align}
 Q_1&<Nx^{-\sigma},                                      \tag{B1}\\
 N^2Q_2Q_3^2&<x^{1-8\sigma},                              \tag{B2}\\
 |a|^\theta NQ_1Q_2^5Q_3^2&<x^{2-15\sigma},               \tag{B3}\\
 N^{1+2\theta}Q_1Q_2^{5-2\theta}Q_3^{2+2\theta}
   &<x^{2-15\sigma}.                                    \tag{B4}
\end{align}
All lengths are at least one. Let $\alpha,\beta,\gamma,\lambda,\nu$
be one-bounded, with the same roughness, full auxiliary SW, and
primitive-modulus support assumptions as Proposition~5.2. All
coefficient and SW bounds are uniform in $a$ if a supremum is taken.
Under Inputs D and S and the stated uniform Lemma~4.7,
\[
 \sum_{q_1\sim Q_1}\sum_{q_2\sim Q_2}\sum_{q_3\sim Q_3}
 \gamma_{q_1}\lambda_{q_2}\nu_{q_3}\Delta(q_1q_2q_3)
 \ll_{A,\sigma} x\log^{-A}x,                             \tag{B5}
\]
uniformly on this residue range. No use is made of the erroneous
deletion error of Lemma~4.6. The displayed $|a|^\theta$ is retained;
it is not replaced by one when applying S.

\paragraph{1. Elementary size consequences.}
Put $n=\log_xN$, $q=\log_xQ_1$, $u=\log_xQ_2$ and
$v=\log_xQ_3$. Dropping the nonnegative $\theta\log_x|a|$
only for the following auxiliary bounds, (B2)--(B3) give
\[
 n+u\le\frac49(2n+u+2v)+\frac19(n+q+5u+2v)
 <\frac23-\frac{47}{9}\sigma.
\]
Consequently
\begin{align}
 Q_2Q_3N^{3/2}&<x^{5/6-119\sigma/18},                    \tag{B6}\\
 Q_1Q_2Q_3&<x^{7/10-13\sigma/2}.                         \tag{B7}
\end{align}
For (B7), use $q<n-\sigma$, half of (B2), and
$u<(2-15\sigma)/5$. Also $N\le x^{1/2-4\sigma}$ and
$N>x^\sigma$ on nonempty length ranges. Thus $M$ is a positive
power of $x$, all relevant moduli are below $x^{1-\delta}$ with
an absolute positive $\delta$, and (B6)--(B7) imply the two size
conditions of Lemma~4.7 with margin $\sigma$.
More precisely, the safe subrange of Input D needed here can use
$\delta=\sigma/3$. For all sufficiently large $x$,
$D=x^{o(1)}\le x^{\sigma/2}$, whence
\[
 N>x^\sigma,\qquad M\ge x^{1/2+4\sigma},\qquad
 QD<Nx^{-\sigma/2},\qquad QD\mathcal R<x^{1-\sigma/3}.
\]
The fixed constants in $D\mathcal R\asymp Q_2Q_3$ are absorbed
by the fixed positive gap in (B7). Thus a preliminary dispersion
repair proved only under these additional power margins is enough;
no extension to arbitrarily short $M$ or $N$ is required here.

\paragraph{2. Smooth parts, including the modulus multiplicity.}
Write
\[
 z_0=x^{1/(\log\log x)^3},\quad y_0=x^{1/\log\log x},
 \quad q_2=u s_2,\quad q_3=v s_3,
\]
where $P^-(u),P^-(v)>z_0$ and $s_2,s_3$ are $z_0$-smooth.
The terms with $s_2>y_0$ or $s_3>y_0$ are negligible with arbitrary
logarithmic saving. Here is a weighted justification of this familiar
smooth-part deletion, rather than multiplying an unweighted estimate
by $\max\tau(q)^B$.

The outer convolution has $|c_b|\le\tau(b)$ and is supported on
$b\asymp x$. Grouping $q_1q_2q_3=l$ gives at most $\tau_3(l)$
representations. The standard primitive progression divisor bound
(Shiu, with zero shift and progression length $\asymp x$) bounds
the absolute counting and mean terms by
$x\tau(l)^C\log^C x/l$, uniformly in the primitive residue.
This uses no moment of $\tau(b-a)$ and no fixed-size representative
for $a$; the residue may first be reduced modulo $l$.
For $\xi=1/\log z_0$, splitting $l$ into its smooth and rough parts,
\[
 \sum_{\substack{l\ll Q_1Q_2Q_3\\\operatorname{sm}(l;z_0)\ge y_0}}
       \frac{\tau(l)^C}{l}
 \ll \log^{C'}x\,y_0^{-\xi}
       \prod_{p\le z_0}\sum_{j\ge0}(j+1)^C p^{-j(1-\xi)}
 \le \exp\{- (\log\log x)^2+O_C(\log\log x)\}.             \tag{B8}
\]
The bound beats every inverse power of $\log x$. The modulus
multiplicity has already been included by increasing $C$.

Dyadically fix the four rough/smooth lengths and the total rough
length $\mathcal R$, with $u\sim U$, $v\sim V$, $uv\sim\mathcal R$,
$\mathcal R\asymp UV$, and $d=s_2s_3\sim D$. There are only
logarithmically many such length choices. Keep all small factors
inside the same dispersion problem; do not sum a logarithmic-error
estimate separately over $x^{o(1)}$ individual values of $d$.
We have $D\ll y_0^2=x^{o(1)}$ and $D\mathcal R\asymp Q_2Q_3$.
Define exactly
\[
 \Lambda_{d,r}=\sum_{s_2s_3=d}
   \sum_{uv=r}\lambda_{us_2}\nu_{vs_3}
       \mathbf1_{\text{all fixed intervals and smooth/rough supports}}.
                                                               \tag{B9}
\]
There is no summation over $q_1$ in this definition and no dependence
on it. Since $(d,r)=1$, $|\Lambda_{d,r}|\le\tau(d)\tau(r)=\tau(dr)$.
The original coefficient $\gamma_{q_1}$ remains outside this convolution.

\paragraph{3. Dispersion and the first exponential sum.}
Apply D with $Q=Q_1$, rough length $\mathcal R$ and the whole dyadic
$d$ range. By (B1), $N>QD\log^Cx$ for all sufficiently large $x$;
the other length hypothesis follows from (B2). The natural Fourier
cutoffs are
\[
 H_1=\frac{QD\mathcal R}{M}\log^5x,
 \qquad H_2=\frac{QD\mathcal R^2}{M}\log^5x.                 \tag{B10}
\]
The coefficient $\Lambda_{d,r}$ is independent of $q$, so Lemma~4.7
applies. Equations (B6)--(B7), and $D\mathcal R\asymp Q_2Q_3$,
give
\[
 |\mathcal E_1|\ll \frac{N^2}{DQ}x^{-\sigma}.              \tag{B11}
\]
Fixed multiplicative constants in the dyadic lengths fit the strict
power margins. The one-bounded inputs $\alpha,\beta$ require no
new truncation before this application. The divisor bound in (B9)
is allowed by D and will be expanded exactly in the second sum.

\paragraph{4. Keep $d$ fixed; do not replace $q$ by $qd$.}
Write the exact second sum as
\[
 \mathcal E_2(a)=\sum_{d\sim D}\frac1d\,\mathcal E_3(a;d), \tag{B12}
\]
where $\mathcal E_3$ has denominator $qr_1r_2$, congruence modulo
$qd$, the two original phases containing $qd$, and cutoff $H_2$.
When bounding a supremum, first take absolute values term by term
in this identity. The source's line2593 inserts an extra $1/d$
in its displayed simplified $\mathcal E'$ although line2564's
$\mathcal E_3$ has none; (B12) fixes the normalization without
relying on that inconsistency.

At fixed $d$, reciprocity gives the exact extra factor
\[
 e\!\left(\frac{ah}{qdr_1r_2n_1}\right).                 \tag{B13}
\]
In the five normalized variables $(q/Q,r_1/\mathcal R,
r_2/\mathcal R,n_1/N,h/H_2)$, its coefficient is
\[
 \frac{aH_2}{Qd\mathcal R^2N}
 =\frac ax\frac Dd\log^5x.
\]
Since $d\sim D$, no extra positive power from $d$ appears in the
variation. The retained-phase candidate, with positive and negative
$h$ split, gives a factor
$x^{5\kappa+o(1)}$, adds one $n_1$ endpoint, and partitions
$(n_1r_2,n_2)$ into the fixed cross-coprime sets.

The outer variable is still the original unrestricted interval of
$q$. Replacing it by $q'=qd$ and then restricting a theorem for all
$q'$ to the multiples of $d$ would be invalid. It is unnecessary:
Lemma~5.1 itself has a fixed $d$, congruence modulo $qd$, and $dq$
inside its inverse phase. Its denominator is $qr_1r_2s_1s_2$,
with no factor $d$, just as required by $\mathcal E_3$.
The Fourier length in (B10) is admissible because
$D\log^5x=x^{o(1)}$ in the lemma's condition
$H'\le x^{o(1)}QR^2S^2/M$.

\paragraph{5. Both small-factor choices must be expanded.}
In $\Lambda_{d,r_1}\overline{\Lambda_{d,r_2}}$, choose independently
$(s_2,s_3)$ and $(\widetilde s_2,\widetilde s_3)$, both with product
$d$. There are at most $\tau(d)^2=x^{o(1)}$ choices. For each,
the rough factors on the two sides have independent one-bounded
sequences
\[
 \lambda_{us_2},\quad\lambda_{u\widetilde s_2},\qquad
 \nu_{vs_3},\quad\nu_{v\widetilde s_3}.
\]
Their actual support intervals, and all original $us_2\sim Q_2$
and $vs_3\sim Q_3$ restrictions, are included in those sequences.
The endpoints for each total rough product are exactly the allowed
cutoffs $r_is_i\le B_i$ of Lemma~5.1.

Apply the non-Hermitian four-phase polarization separately in the
two $\lambda$ sequences and the two $\nu$ sequences. After dividing
each sum of two one-bounded sequences by2, this costs at most
$4^2=16$. The independent $n_1$ endpoint costs at most another4.
Thus the total constant is at most64, besides $\tau(d)^2$; no equality
of the small-factor choices is assumed. Polarization is after the
dispersion use of SW, so the resulting sequences need only be bounded.
The bounded one-sided $h$ weights are permitted by the proof-level
extension of S; they are not arbitrary $q$ weights.

In Lemma~5.1 notation use $S=U=Q_2'$ and $R=V=Q_3'$.
All powers of $Q_2',Q_3'$ in (B2)--(B4) are positive, including
$5-2\theta$, so their size conditions with $7\sigma,14\sigma$
follow from (B2)--(B4), after the fixed dyadic constants. Therefore
\[
 |\mathcal E_3(a;d)|
 \ll\frac{N^2}{Q}x^{-\sigma+5\kappa+o(1)}.                \tag{B14}
\]
The subpower term includes small-factor pairs, coloring, divisor
normalizations, product endpoints and logarithmic variation.
It is uniform for $d\ll y_0^2$ and the fixed SW family.

\paragraph{6. Final $d$ sum and the square root.}
In a dyadic $d$ interval, $\sum_{d\sim D}1/d=O(1)$.
Equation (B12) consequently preserves (B14) without an artificial
factor $D$. To match the normalization of D, absorb the additional
$D=x^{o(1)}$ as well. Choose $x$ sufficiently large that the whole
subpower loss in these steps is at most $x^{\sigma/10}$. Since
$5\kappa\le\sigma/20$, this proves the deliberately weaker bound
\[
 |\mathcal E_1|+|\mathcal E_2|
 \ll\frac{N^2}{DQ}x^{-\sigma/2}.                          \tag{B15}
\]
Substitution into (D) yields
\[
 M(DQ)^{1/2}(|\mathcal E_1|^{1/2}+|\mathcal E_2|^{1/2})
 \ll x^{1-\sigma/4}.
\]
All remaining fixed logarithmic powers are harmless, for example
giving $O(x^{1-\sigma/5})$ for sufficiently large $x$. Take the
logarithmic exponent in D larger than the requested $A$ plus the
finite number of dyadic decompositions. Together with (B8) this
proves (B5), subject exactly to the inputs at the beginning.

\paragraph{7. Coefficient scope and the unresolved earlier boundary.}
This proposition is the one-bounded Proposition~5.2 bridge. It is
not automatically a theorem for arbitrary divisor-bounded HB
convolutions. A log-height truncation of a rough SW sequence preserves
SW by the separate Lemma~4.3, but discarding the large coefficients
in moving progressions requires a valid uniform shifted-divisor
estimate and handling of $nm=a$. No $x^{o(1)}$ multiplier is allowed
to be silently applied to a merely logarithmic remainder. In this
proof, the only $x^{o(1)}$ multiplicities occur after dispersion and
are paid by a power saving; the original outer coefficients are
already one-bounded.

Before the separate supplements assembled in section~8, the earliest
unjustified statement in a direct invocation of the source was Input D
for the moving residue. Maynard~I
Lemmas14.2--14.3 use $\tau(mn-a)$ and a fixed-shift divisor-moment
bound. The formal expressions at $mn=a$ contain $\tau(0)$ and
cannot be justified by the fixed-residue text. In particular the
extra condition $Q_1Q_2Q_3<\sqrt{xN}$ does not follow from
(B1)--(B4); it may not be used to hide this issue. A separate
zero-deletion/shifted-moment reconstruction is necessary to supply D.
The present note does not claim it by fiat. Nor does this note
re-prove the spectral estimate in S or the downstream small-product
tail and finite-tree applications. It provides the complete
$D+S\Longrightarrow$(B5) bridge with the disputed deletion step
removed and all small-factor and endpoint losses exposed.

\paragraph{8. Restricted assembly with the separately supplied repairs.}
The preceding proposition separates its inputs so that it also records
precisely what remains conditional in the general $\theta$, all-$a$
form.  At completion of this note three new, separately reviewed
documents supply the needed inputs in the following smaller subclass:
\[
 \theta=7/32,\qquad c_0x\le |a|\le x^{1+\kappa},\qquad
 0\le\kappa\le\sigma/100,
\]
where $c_0>0$ is fixed and all four conditions (B1)--(B4) are retained.
The zero-product dispersion supplement and its independent review
supply D for $\delta=\sigma/3$.  Section~1 proves every extra power
margin that supplement requires.  The large-residue Lemma~5.1 audit
supplies S for the actual coprime spectral subclass: its level is
$z=s_1s_2$ with $(a,z)=1$, a condition already present and preserved
in the inner rough-factor sums here.  It replaces the printed
phase-free same-eigenvector shortcut by a packet norm with retained
phase; no unrestricted non-coprime spectral assertion is needed.

Finally the direct Maynard~I Lemma~17.2 supplement, with the independent
E1 review, gives (B11) with the weaker exponent $\sigma/2$, which is
exactly enough for (B15).  It uses $|a|\ge c_0x$ and $\theta=7/32$;
these hypotheses may not be dropped.  Its allowed $q$-dependent
coefficients justify the fixed-$d$ substitution $q'=dq$ in E1 only.
This is different from the E2 argument of section~4.

Consequently (B5) is derived for this restricted subclass, relative
to the explicitly named classical Shiu, Weil, Poisson, SW and spectral
inputs in those three documents.  The composition is an analytic
derivation, not a machine proof or a reproof of those external
theorems.  It does not assert the original complete large-residue
Theorem~1.7 or the unrestricted version of Theorem~1.8.  The original
deletion proof remains invalid as written.  Project-wide HB and
small-product applications and the separate reported Wu issue are
outside this proposition; no global numerical record is reinstated.
